% ==============================================================================
% Section 3.4: Cross-Method Comparison and Unified Architecture (Clean Skeleton)
% Owner: Member 3 (TV3) | Target length: ~2 pages
% ==============================================================================

\section{Tổng hợp mối liên hệ, độ phức tạp và khả năng kết hợp}
\label{sec:method_integration}

\draftnote[TV3]{Mục 3.4 dài khoảng 2 trang. Đây là mục tổng hợp liên kết toàn bộ 4 phương pháp kỹ thuật theo khía cạnh: độ phức tạp, phạm vi giải quyết, khả năng ghép nối thành một kiến trúc đầu--cuối và các điểm nghẽn khi tích hợp.}

\subsection{Ma trận so sánh toàn diện bốn phương pháp}
\label{subsec:4method_matrix}

\draftnote[TV3]{Hoàn thiện Bảng~\ref{tab:four_methods_summary} đối chiếu UniversalRAG, ColPali, VisRAG và CoRe-MMRAG theo các trục: Tầng xử lý, Bài toán giải quyết, Chi phí/Độ phức tạp (thời gian, bộ nhớ) và Điều kiện thất bại.}

\begin{table}[H]
\centering
\small
\caption[So sánh tổng hợp 4 phương pháp trọng tâm]{So sánh tổng hợp 4 phương pháp trọng tâm theo tầng xử lý, phạm vi và độ phức tạp.}
\label{tab:four_methods_summary}
\renewcommand{\arraystretch}{1.24}
\begin{tabularx}{\textwidth}{@{}P{3.0cm}P{2.3cm}L L P{2.7cm}@{}}
\toprule
\rowcolor{tblhead}
\thead{Phương pháp} & \thead{Tầng xử lý} & \thead{Bài toán giải quyết} & \thead{Chi phí / Độ phức tạp} & \thead{Điều kiện thất bại} \\
\midrule
UniversalRAG & \dots & \dots & \dots & \dots \\
\rowcolor{zebrabg}
ColPali & \dots & \dots & \dots & \dots \\
VisRAG & \dots & \dots & \dots & \dots \\
\rowcolor{zebrabg}
CoRe-MMRAG & \dots & \dots & \dots & \dots \\
\bottomrule
\end{tabularx}
\end{table}

\subsection{Kiến trúc tích hợp đề xuất và các điểm nghẽn khi kết hợp}
\label{subsec:unified_architecture}

\draftnote[TV3]{Đề xuất kịch bản ghép nối cả 4 phương pháp thành một pipeline tham chiếu thống nhất (UniversalRAG định tuyến $\rightarrow$ ColPali/VisRAG xử lý kho tài liệu trực quan $\rightarrow$ CoRe-MMRAG hòa giải xung đột PRKI/VTKI $\rightarrow$ tạo sinh câu trả lời). Có thể đặt nhận định cốt lõi vào môi trường \texttt{keybox}. Phân tích 3 điểm nghẽn thực tế khi tích hợp: cộng dồn độ trễ, lan truyền sai số xuyên tầng (cascading errors) và chi phí đồng bộ nhiều mô hình nền khác nhau.}
